% sec1_1_part1.tex — Section 1.1 (part 1): Linearity, Matrix Notation, Strict Exogeneity

In this section we present the assumptions that comprise the classical linear regression
model.  In the model, the variable in question (called the \textbf{dependent variable},
the \textbf{regressand}, or more generically the \textbf{left-hand[-side] variable}) is
related to several other variables (called the \textbf{regressors}, the
\textbf{explanatory variables}, or the \textbf{right-hand[-side] variables}).  Suppose we
observe $n$ values for those variables.  Let $y_i$ be the $i$-th observation of the
dependent variable in question and let $(x_{i1}, x_{i2}, \ldots, x_{iK})$ be the $i$-th
observation of the $K$ regressors.  The \textbf{sample} or \textbf{data} is a collection
of those $n$ observations.

The data in economics cannot be generated by experiments (except in experimental
economics), so both the dependent and independent variables have to be treated as random
variables, variables whose values are subject to chance.  A \textbf{model} is a set of
restrictions on the joint distribution of the dependent and independent variables.  That
is, a model is a set of joint distributions satisfying a set of assumptions.  The
classical regression model is a set of joint distributions satisfying Assumptions
1.1--1.4 stated below.

%% ─── The Linearity Assumption ───────────────────────────────────────────────
\subsection*{The Linearity Assumption}

The first assumption is that the relationship between the dependent variable and the
regressors is linear.

\begin{assumption}[linearity]
\label{ass:linearity}
\begin{equation}
  y_i = \beta_1 x_{i1} + \beta_2 x_{i2} + \cdots + \beta_K x_{iK} + \varepsilon_i
  \quad (i = 1, 2, \ldots, n),
  \label{eq:1.1.1}
\end{equation}
\emph{where $\beta$'s are unknown parameters to be estimated, and $\varepsilon_i$ is the
unobserved error term with certain properties to be specified below.}
\end{assumption}

The part of the right-hand side involving the regressors,
$\beta_1 x_{i1} + \beta_2 x_{i2} + \cdots + \beta_K x_{iK}$, is called the
\textbf{regression} or the \textbf{regression function}, and the coefficients ($\beta$'s)
are called the \textbf{regression coefficients}.  They represent the marginal and separate
effects of the regressors.  For example, $\beta_2$ represents the change in the dependent
variable when the second regressor increases by one unit while other regressors are held
constant.  In the language of calculus, this can be expressed as
$\partial y_i / \partial x_{i2} = \beta_2$.  The linearity implies that the marginal
effect does not depend on the level of regressors.  The error term represents the part of
the dependent variable left unexplained by the regressors.

\begin{example}[consumption function]
\label{ex:consumption}
The simple consumption function familiar from introductory economics is
\begin{equation}
  CON_i = \beta_1 + \beta_2 YD_i + \varepsilon_i,
  \label{eq:1.1.2}
\end{equation}
where $CON$ is consumption and $YD$ is disposable income.  If the data are annual
aggregate time-series, $CON_i$ and $YD_i$ are aggregate consumption and disposable income
for year~$i$.  If the data come from a survey of individual households, $CON_i$ is
consumption by the $i$-th household in the cross-section sample of $n$ households.  The
consumption function can be written as~\eqref{eq:1.1.1} by setting $y_i = CON_i$,
$x_{i1} = 1$ (a constant), and $x_{i2} = YD_i$.  The error term $\varepsilon_i$
represents other variables besides disposable income that influence consumption.  They
include those variables---such as financial assets---that might be observable but the
researcher decided not to include as regressors, as well as those variables---such as the
``mood'' of the consumer---that are hard to measure.  When the equation has only one
nonconstant regressor, as here, it is called the \textbf{simple regression model}.
\end{example}

The linearity assumption is not as restrictive as it might first seem, because the
dependent variable and the regressors can be transformations of the variables in question.
Consider

\begin{example}[wage equation]
\label{ex:wage}
A simplified version of the wage equation routinely estimated in labor economics is
\begin{equation}
  \log(WAGE_i) = \beta_1 + \beta_2 S_i + \beta_3 TENURE_i + \beta_4 EXPR_i
    + \varepsilon_i,
  \label{eq:1.1.3}
\end{equation}
where $WAGE ={}$ the wage rate for the individual, $S ={}$ education in years,
$TENURE ={}$ years on the current job, and $EXPR ={}$ experience in the labor force
(i.e., total number of years to date on all the jobs held currently or previously by the
individual).  The wage equation fits the generic format~\eqref{eq:1.1.1} with
$y_i = \log(WAGE_i)$.  The equation is said to be in the \textbf{semi-log} form because
only the dependent variable is in logs.  The equation is derived from the following
nonlinear relationship between the level of the wage rate and the regressors:
\begin{equation}
  WAGE_i = \exp(\beta_1)\,\exp(\beta_2 S_i)\,\exp(\beta_3 TENURE_i)\,
    \exp(\beta_4 EXPR_i)\,\exp(\varepsilon_i).
  \label{eq:1.1.4}
\end{equation}
By taking logs of both sides of~\eqref{eq:1.1.4} and noting that
$\log[\exp(x)] = x$, one obtains~\eqref{eq:1.1.3}.  The coefficients in the semi-log
form have the interpretation of \emph{percentage changes}, not changes in levels.  For
example, a value of $0.05$ for $\beta_2$ implies that an additional year of education has
the effect of raising the wage rate by 5 percent.  The difference in the interpretation
comes about because the dependent variable is the log wage rate, not the wage rate itself,
and the change in logs equals the percentage change in levels.
\end{example}

Certain other forms of nonlinearities can also be accommodated.  Suppose, for example,
the marginal effect of education tapers off as the level of education gets higher.  This
can be captured by including in the wage equation the squared term $S^2$ as an additional
regressor.  If the coefficient of the squared term is $\beta_5$, the marginal effect of
education is
\[
  \beta_2 + 2\beta_5 S \;(= \partial\log(WAGE)/\partial S).
\]
If $\beta_5$ is negative, the marginal effect of education declines with the level of
education.

There are, of course, cases of genuine nonlinearity.  For example, the
relationship~\eqref{eq:1.1.4} could not have been made linear if the error term entered
additively rather than multiplicatively:
\[
  WAGE_i = \exp(\beta_1)\,\exp(\beta_2 S_i)\,\exp(\beta_3 TENURE_i)\,
    \exp(\beta_4 EXPR_i) + \varepsilon_i.
\]
Estimation of nonlinear regression equations such as this will be discussed in Chapter~7.

%% ─── Matrix Notation ────────────────────────────────────────────────────────
\subsection*{Matrix Notation}

Before stating other assumptions of the classical model, we introduce the vector and
matrix notation.  The notation will prove useful for stating other assumptions precisely
and also for deriving the OLS estimator of~$\bm{\beta}$.  Define $K$-dimensional (column)
vectors $\vec{x}_i$ and $\bm{\beta}$ as
\begin{equation}
  \underset{(K \times 1)}{\vec{x}_i}
  = \begin{bmatrix} x_{i1} \\ x_{i2} \\ \vdots \\ x_{iK} \end{bmatrix}, \quad
  \underset{(K \times 1)}{\bm{\beta}}
  = \begin{bmatrix} \beta_1 \\ \beta_2 \\ \vdots \\ \beta_K \end{bmatrix}.
  \label{eq:1.1.5}
\end{equation}
By the definition of vector inner products,
$\vec{x}_i'\bm{\beta} = \beta_1 x_{i1} + \beta_2 x_{i2} + \cdots + \beta_K x_{iK}$.
So the equations in Assumption~1.1 can be written as
\begin{equation}
  y_i = \vec{x}_i'\bm{\beta} + \varepsilon_i \quad (i = 1, 2, \ldots, n).
  \tag{\ref{eq:1.1.1}$'$}
  \label{eq:1.1.1p}
\end{equation}
Also define
\begin{equation}
  \underset{(n \times 1)}{\vec{y}}
  = \begin{bmatrix} y_1 \\ \vdots \\ y_n \end{bmatrix}, \quad
  \underset{(n \times 1)}{\bm{\varepsilon}}
  = \begin{bmatrix} \varepsilon_1 \\ \vdots \\ \varepsilon_n \end{bmatrix}, \quad
  \underset{(n \times K)}{\vec{X}}
  = \begin{bmatrix} \vec{x}_1' \\ \vdots \\ \vec{x}_n' \end{bmatrix}
  = \begin{bmatrix}
      x_{11} & \cdots & x_{1K} \\
      \vdots & \cdots & \vdots \\
      x_{n1} & \cdots & x_{nK}
    \end{bmatrix}.
  \label{eq:1.1.6}
\end{equation}
In the vectors and matrices in~\eqref{eq:1.1.6}, there are as many rows as there are
observations, with the rows corresponding to the observations.  For this reason $\vec{y}$
and $\vec{X}$ are sometimes called the \textbf{data vector} and the \textbf{data matrix}.
Since the number of columns of $\vec{X}$ equals the number of rows of $\bm{\beta}$,
$\vec{X}$ and $\bm{\beta}$ are conformable and $\vec{X}\bm{\beta}$ is an $n \times 1$
vector.  Its $i$-th element is $\vec{x}_i'\bm{\beta}$.  Therefore, Assumption~1.1 can be
written compactly as
\[
  \underset{(n \times 1)}{\vec{y}}
  = \underbrace{
      \underset{(n \times K)}{\vec{X}} \;
      \underset{(K \times 1)}{\bm{\beta}}
    }_{(n \times 1)}
  + \underset{(n \times 1)}{\bm{\varepsilon}}.
\]

%% ─── The Strict Exogeneity Assumption ───────────────────────────────────────
\subsection*{The Strict Exogeneity Assumption}

The next assumption of the classical regression model is

\begin{assumption}[strict exogeneity]
\label{ass:strict_exogeneity}
\begin{equation}
  \E(\varepsilon_i \mid \vec{X}) = 0 \quad (i = 1, 2, \ldots, n).
  \label{eq:1.1.7}
\end{equation}
\end{assumption}

Here, the expectation (mean) is conditional on the regressors for \emph{all} observations.
This point may be made more apparent by writing the assumption without using the data
matrix as
\[
  \E(\varepsilon_i \mid \vec{x}_1, \ldots, \vec{x}_n) = 0
  \quad (i = 1, 2, \ldots, n).
\]
To state the assumption differently, take, for any given observation~$i$, the joint
distribution of the $nK + 1$ random variables,
$f(\varepsilon_i, \vec{x}_1, \ldots, \vec{x}_n)$, and consider the conditional
distribution, $f(\varepsilon_i \mid \vec{x}_1, \ldots, \vec{x}_n)$.  The conditional mean
$\E(\varepsilon_i \mid \vec{x}_1, \ldots, \vec{x}_n)$ is in general a nonlinear function
of $(\vec{x}_1, \ldots, \vec{x}_n)$.  The strict exogeneity assumption says that this
function is a constant of value zero.\footnote{Some authors define the term ``strict
exogeneity'' somewhat differently.  For example, in Koopmans and Hood (1953) and Engle,
Hendry, and Richards (1983), the regressors are strictly exogenous if $\vec{x}_i$ is
independent of $\varepsilon_j$ for all $i$, $j$.  This definition is stronger than, but
not inconsistent with, our definition of strict exogeneity.}

Assuming this constant to be zero is not restrictive if the regressors include a constant,
because the equation can be rewritten so that the conditional mean of the error term is
zero.  To see this, suppose that $\E(\varepsilon_i \mid \vec{X})$ is $\mu$ and
$x_{i1} = 1$.  The equation can be written as
\begin{align*}
  y_i &= \beta_1 + \beta_2 x_{i2} + \cdots + \beta_K x_{iK} + \varepsilon_i \\
      &= (\beta_1 + \mu) + \beta_2 x_{i2} + \cdots + \beta_K x_{iK}
         + (\varepsilon_i - \mu).
\end{align*}
If we redefine $\beta_1$ to be $\beta_1 + \mu$ and $\varepsilon_i$ to be
$\varepsilon_i - \mu$, the conditional mean of the new error term is zero.  In virtually
all applications, the regressors include a constant term.
